% Generated draft: requires mathematical and source review.
\documentclass[11pt]{article}
\usepackage[margin=25mm]{geometry}
\usepackage{fontspec}
\setmainfont{Latin Modern Roman}
\usepackage{amsmath,amssymb,mathtools,unicode-math}
\setmathfont{Latin Modern Math}
\usepackage{enumitem,needspace,xurl,hyperref}
\setlength{\parindent}{0pt}
\setlength{\parskip}{5pt}
\setlength{\emergencystretch}{3em}
\newcommand{\N}{\mathbb{N}}
\newcommand{\Z}{\mathbb{Z}}
\newcommand{\Q}{\mathbb{Q}}
\newcommand{\R}{\mathbb{R}}
\newcommand{\C}{\mathbb{C}}
\newcommand{\F}{\mathbb{F}}
\newcommand{\A}{\mathbb{A}}
\newcommand{\PP}{\mathbb{P}}
\newcommand{\E}{\mathbb{E}}
\newcommand{\eps}{\varepsilon}
\newcommand{\dd}{\,\mathrm{d}}
\begin{document}
\section*{UnsolvedMath ID: 30001625}
\textbf{Source problem code:} catalogue entry
\subsection{Definitions and mathematical statement}
For each positive integer \(n\), put
\[
R_n=\C[t_1^{\pm1},\ldots,t_n^{\pm1}],
\]
and let \(\mathfrak g\) be a finite-dimensional semisimple complex Lie algebra. Here the nullity parameter is the number \(n\) of Laurent variables.

An \(R_n\)-form of \(\mathfrak g\otimes_{\C}R_n\) means an \(R_n\)-Lie algebra \(\mathcal L\) for which there exists a faithfully flat commutative \(R_n\)-algebra \(S\) and an isomorphism of \(S\)-Lie algebras
\[
\mathcal L\otimes_{R_n}S\simeq\mathfrak g\otimes_{\C}S.
\]
The classification below is interpreted as classification up to \(R_n\)-linear Lie algebra isomorphism. The supplied record does not explicitly specify the descent topology or equivalence relation; this is a conventional interpretation requiring review against the source.

To define a multiloop algebra, choose positive integers \(m_1,\ldots,m_n\), primitive \(m_i\)-th roots of unity \(\zeta_i\in\C\), and pairwise commuting automorphisms \(\sigma_i\in\operatorname{Aut}_{\C}(\mathfrak g)\) satisfying \(\sigma_i^{m_i}=1\). For \(\mathbf k=(k_1,\ldots,k_n)\in\Z^n\), set
\[
\mathfrak g_{\overline{\mathbf k}}
=\{x\in\mathfrak g: \sigma_i(x)=\zeta_i^{k_i}x\text{ for every }i\}.
\]
The corresponding multiloop algebra is
\[
M(\mathfrak g;\boldsymbol\sigma,\mathbf m)
=\bigoplus_{\mathbf k\in\Z^n}
\mathfrak g_{\overline{\mathbf k}}\otimes_{\C}
\C z_1^{k_1}\cdots z_n^{k_n}
\subseteq\mathfrak g\otimes_{\C}\C[z_1^{\pm1},\ldots,z_n^{\pm1}],
\]
with bracket inherited from the ambient Lie algebra and \(R_n\)-action specified by \(t_i=z_i^{m_i}\). In this draft, a form is called multiloop if it is \(R_n\)-isomorphic to such an algebra.

\textbf{Classification problem.} For every positive integer \(n\) and every finite-dimensional complex semisimple \(\mathfrak g\), classify all \(R_n\)-forms of \(\mathfrak g\otimes_{\C}R_n\), and characterize precisely those forms that are multiloop algebras.

\textbf{Nullity-two question.} For \(n=2\), the supplied record additionally asks whether every form that is not multiloop is of ``Margaux type.'' No definition of ``Margaux type'' or ``Margaux-like'' is supplied, so this additional question cannot yet be expressed as a precise quantified mathematical assertion. Moreover, the original wording refers to counterexamples to being ``loop algebras,'' whereas the cleaned wording refers to non-multiloop forms. Their intended agreement requires source review; no equivalence of those terms is asserted here.
\subsection{Short English statement}
Classify forms of a finite-dimensional complex semisimple Lie algebra over a Laurent polynomial ring, and determine which arise from commuting finite-order automorphisms by the multiloop construction. In two variables, the source also asks whether all exceptions are of Margaux type, but that term and the original use of ``loop'' need clarification. This is a statement-only draft requiring human mathematical review.
\subsection{Sources}
\begin{itemize}
\item Local UnsolvedMath registry, \texttt{data/unsolvedmath/problems.json}, permanent ID \(30001625\), source problem code \texttt{catalogue entry}, title ``Forms and Multiloop Algebras in Nullity Two,'' in the supplied Oberwolfach Reports Open Problems collection. No individual proposer is identified in the supplied record; the proposed year is recorded as 2010.
\item Supplied source citation: \emph{Infinite Dimensional Lie Theory} (2011), \emph{Oberwolfach Reports}, pp.~2977--3042, DOI \texttt{10.4171/owr/2010/51}. Original OWR report: \url{https://doi.org/10.4171/owr/2010/51}. The supplied background attributes the problem to workshop open-problem discussions documented in this report.
\item Supplied literature reference: Allison--Berman--Pianzola, \url{https://arxiv.org/abs/1002.2674}.
\item Supplied literature reference: Gille--Pianzola, \url{https://arxiv.org/abs/1109.3405}.
\item Literature assessment dated exactly \texttt{2026-08-21}: the supplied curation labels the problem ``partially solved'' and reports that nullity-one forms are classified as loop algebras, that nullity-two multiloop algebras and several classical cases are classified, and that Margaux-type counterexamples leave the full classification of all \(R_n\)-forms unresolved. These are supplied dated curation claims, not independently verified mathematical or current-status conclusions. The supplied background explicitly requests independent verification of time-sensitive status and cited claims.
\item The record labels its statement \texttt{corrected\_verified} and reports that it was checked against the official Oberwolfach report. This verification attribution is supplied metadata; no source was independently visited for this draft.
\end{itemize}
\end{document}
